% =============================================================================
% SHORT WORKSHOP NOTE (portfolio §3.3) — demoted rebuild of P1
% Zero GPU. Claim fixes already in paper_P1_scaling; this is the shippable short unit.
% Build: latexmk -pdf paper_P1_identifiability_note.tex
% =============================================================================
\documentclass[11pt]{article}

\usepackage[utf8]{inputenc}
\usepackage[T1]{fontenc}
\usepackage[margin=1in]{geometry}
\usepackage{hyperref}
\usepackage{booktabs}
\usepackage{array}
\usepackage{amsmath,amssymb}
\usepackage{microtype}

\title{Failed Identifiability of GRPO Post-Training Scaling\\
\large A Short Negative Note on Cross-Scale Anchors}
\author{Arvind C R\\
PES University\\\texttt{arvindcr4@gmail.com}}
\date{\today}

\begin{document}
\maketitle

\begin{abstract}
We audit whether the available group-relative policy optimization (GRPO)
post-training traces support a size-scaling law. They do not currently support
one: short, heterogeneous horizons, ceiling effects, and weakly identified
curve fits prevent a positive scaling interpretation. This is an evidence
limitation, not a statistical demonstration of zero slope or a universal
failure of scaling. Earlier permutation, model-selection, architecture-gap,
and leave-one-anchor claims are not retained as inferential conclusions.
A Nemotron-3-Super-120B-A12B trace contains 11 zero-reward steps out of 20;
this is a descriptive run observation, not a model-family comparison.
\end{abstract}

\section{What is (and is not) claimed}

\paragraph{Claim.} The currently documented anchors are insufficient to
identify a positive GRPO size trend or a universal saturation exponent.

\paragraph{Non-claims.} No multi-seed powered cross-scale significance test; no
causal architecture ranking; no claim that larger models cannot improve under
other stacks, horizons, or tasks. Frontier anchors are single-seed and
descriptive by design. Predictions recorded in an internal iteration ledger are
\emph{not} an immutable external preregistration.

\section{Panel and provenance limits}

The historical 12-row fit panel is explicitly enumerated in
\nolinkurl{platform_hybrid/experiments/results/scaling_law_iter109_saturation.tsv}:
Qwen3.5-4B (30 steps), Qwen3-8B (30), Llama-3.1-8B-Instruct (30),
Qwen3-32B (3), Qwen3.5-27B (3), gpt-oss-20B (30), Qwen3-30B-MoE (5),
Qwen3-30B-MoE-Inst (3), DeepSeek-V3.1 (20),
Nemotron-3-Super-120B-A12B (20; historically labelled Nemotron-120B),
Qwen3-235B-MoE (4), and Kimi-K2-Thinking (20).
Its \texttt{trace\_file} column identifies the input filename for each row.
The table's \texttt{params\_B} is nominal total parameter count, not active
parameters or per-token compute; its Nemotron \texttt{dense} classification
is wrong (LatentMoE, 120B total / 12B active).

These rows do not constitute a controlled 0.6B--1T experiment. A complete
row-level mapping of checkpoint revisions, prompt IDs/splits, seeds,
backend versions, decoding, and stopping rules is not supplied by that fit
TSV. In particular, ceiling-saturated instruct models and short probes cannot
be treated as comparable completed optimization trajectories. We therefore
withdraw panel-wide inferential conclusions pending that provenance audit.
The source repository is \url{https://github.com/arvindcr4/tinker-rl-lab};
this note does not provide a frozen archival release.

\section{Descriptive findings and withdrawn analyses}

\begin{table}[ht]
\centering
\caption{Audit status. Historical analyses are identified, not endorsed as tests.}
\label{tab:neg}
\small
\begin{tabular}{@{}p{0.42\linewidth}>{\raggedright\arraybackslash}p{0.54\linewidth}@{}}
\toprule
Check & Status \\
\midrule
Size regressions on fitted rates and asymptotes &
Historical $p=0.7382$ and $0.5061$ in
\nolinkurl{scaling_law_iter109b_permtest.tsv}; not retained as evidence of
flatness because fitted boundary values and heterogeneous traces are not
exchangeable scientific replicates. \\
Saturation-rate fits &
Historical form $R(t)=R_\infty+(R_0-R_\infty)e^{-\lambda t}$ in
\nolinkurl{platform_modal/scripts/scaling_law_iter109.py}, with bounds
$R_0\in[-1.5,1.5]$, $R_\infty\in[0,1.5]$, $\lambda\in[10^{-4},50]$.
These bounds permit nonphysical reward asymptotes. No universal exponent is inferred. \\
Corrected Akaike information criterion (AICc) comparisons &
Withdrawn as headline evidence; short per-family panels and uncertain fitted
estimands do not support the former blanket constant-model preference. \\
Nemotron-3-Super-120B-A12B &
Descriptive trace: peak reward $0.875$ at step 3; last-10 $0.1625$;
zero-reward steps $11/20$; $5/20$ steps above $0.25$.
Input: \nolinkurl{frontier_gsm8k_nemotron-120b.json}. \\
MoE--dense comparison &
Withdrawn: historical split misclassified Nemotron. Neither the original
nor the corrected small-panel $p$-value identifies an architecture effect. \\
Leave-Nemotron analysis &
Withdrawn: no trusted sensitivity receipt is established here; earlier
claims of robustness or reversal are not retained. \\
\bottomrule
\end{tabular}
\end{table}

\section{Interpretation}

Plateau, drift, and collapse are descriptive trajectory labels, not an
established stack-conditioned taxonomy. Recipe, horizon, prompt difficulty,
checkpoint capability and size co-vary. This audit therefore supports neither
a universal positive scaling law nor its universal negation. Repeated-seed,
matched-budget experiments with documented checkpoints and held-out prompts
are required before comparing architecture or scaling slopes. No standalone
scientific novelty for these generic identifiability limitations is claimed.

\section{Document scope}
This correction note withdraws overstatements in the historical scaling audit.
It does not certify the longer draft or replace missing panel provenance.

\end{document}
